%% kapa: exam=1 round=33 year=2022 period=1 form=A subject=civil_law subject_ko=민법 seq=25 number=25
%% source: Q-Net 감정평가사 기출문제 (cst00309) artlSeq=5209997 file=raw/1008/5209997/2022년도 제33회 감정평가사 1차시험 1교시(민법,경제학,부동산학).pdf
%% conversion: kordoc 4.24.51 -> tools/gen1.py (auto); review=proofread-sonnet; flags=none
\begin{question}{25}
선의취득에 관한 설명으로 옳지 않은 것은? (다툼이 있으면 판례에 따름)
\begin{choices}
  \item 경매에 의해서는 동산을 선의취득할 수 없다.
  \item 점유개정에 의한 인도로는 선의취득이 인정되지 않는다.
  \item 동산질권도 선의취득할 수 있다.
  \item 선의취득자는 임의로 선의취득의 효과를 거부하고 종전 소유자에게 동산을 반환받아 갈 것을 요구할 수 없다.
  \item 점유보조자가 횡령한 물건은 민법 제250조의 도품·유실물에 해당하지 않는다.
\end{choices}
\end{question}
\answer{1}
